\documentclass[11 pt , a 4 paper ]{article}
\usepackage[dvipsnames]{xcolor}
\usepackage{tikz}
\usetikzlibrary{shapes}
\usetikzlibrary{calc}
\usepackage{anyfontsize}
\usepackage{sectsty}
\usepackage{graphicx} % Required for inserting images
\usepackage{amsmath} % AMS Base Package
\usepackage{amssymb} % AMS Symbols
\usepackage{amsthm} % AMS
\usepackage{hyperref} % Hyperrefferings
\usepackage{amsfonts} % AMS Fonts
\usepackage{tabto}
\usepackage{enumitem}
\usepackage[utf8]{inputenc}
\usepackage{pgfplots} % Plots
\pgfplotsset{width=10cm,compat=1.9}
\usepackage[margin=1cm]{geometry}
\usepackage[skins]{tcolorbox}
\usepackage{listings}  %Code blocks
\usepackage{enumitem} % Fancy Bullet points
\usepackage{array} %Tables
\usepackage{xcolor} % Text Color
\usepackage{svg}
\usepackage{pdfpages}
\usepackage{array} %Tables
\usepackage[table]{xcolor} % Text Color
\usepackage{float}
\usepackage{booktabs} % For professional horizontal lines
\usepackage{siunitx}  % For aligning numbers on the decimal point
\usepackage[table]{xcolor} % For adding color (e.g., zebra stripes)
\usepackage{tabularx} % For tables that adjust to a specific width
\usepackage{ragged2e} % for '\RaggedRight' macro (suppress full justification)
\usepackage{booktabs} % for '\toprule', '\midrule', and '\bottomrule' macros

\definecolor{aqua}{RGB}{3,164,186}
\definecolor{grey}{RGB}{200,200,200}


{\Huge\title {UC DEVS SHOWCASE}}

\author{Ruhan Shafi}
\date{\today}

\theoremstyle{definition}
\newtheorem{thm}{Defination}


\usepackage{xcolor}

\definecolor{codegreen}{rgb}{0,0.6,0}
\definecolor{codegray}{rgb}{0.5,0.5,0.5}
\definecolor{codepurple}{rgb}{0.58,0,0.82}
\definecolor{backcolour}{rgb}{0.95,0.95,0.92}

\lstdefinestyle{mystyle}{
	backgroundcolor=\color{backcolour},   
	commentstyle=\color{codegreen},
	keywordstyle=\color{magenta},
	numberstyle=\tiny\color{codegray},
	stringstyle=\color{codepurple},
	basicstyle=\ttfamily\footnotesize,
	breakatwhitespace=false,         
	breaklines=true,                 
	captionpos=b,                    
	keepspaces=true,                 
	numbers=left,                    
	numbersep=5pt,                  
	showspaces=false,                
	showstringspaces=false,
	showtabs=false,                  
	tabsize=2
}

\lstset{style=mystyle}
\begin{document}
	\ \vskip -1cm
	\rightline{\includegraphics[scale=.06]{~/Notes/LaTex Work/logo3.png}}
	
	\vskip -40pt
	
	\noindent
	{UC DEVS Showcase | Ruhan Shafi (u3284342)}
	
	\vskip 5pt
	
	\noindent{\bf\Large CNN Showcase: Image Classification and Model Explainer}
	
	\vskip 7pt
	{\color{aqua}\hrule height 1pt}
	{\color{grey}\hrule height 1pt}
	
	\vskip 5pt

\section{What is Machine Learning?}

Machine learning is a subfield of artificial intelligence concerned with the development of computational systems that improve their performance on a given task through exposure to data, rather than through explicit, rule-based programming. Instead of a programmer writing down every rule the computer should follow, a machine learning model is shown many examples and gradually learns the underlying patterns for itself.

\section{What is a CNN?}

 
A Convolutional Neural Network (CNN) is a type of deep learning algorithm that is particularly effective for analysing visual data. It is designed to automatically and adaptively learn spatial hierarchies of features from input images, meaning it starts by noticing very simple patterns and builds up to recognising complex objects. CNNs are widely used in image classification, object detection, and other computer vision tasks, and are the type of model used in this showcase.

\section{How does a CNN work?}

\begin{itemize}[leftmargin=1.5em]
	\item \textbf{Convolutional layers} slide a small grid of numbers, called a \emph{filter}, across the image. Each filter is tuned to react strongly to a particular simple pattern, such as a vertical edge, a curve, or a patch of a certain colour. The result of sliding a filter across the whole image is called a \emph{feature map}.
	\item \textbf{Pooling layers} shrink the feature maps down, keeping only the most important information from small regions. This makes the network faster to run and less sensitive to a feature being a few pixels out of place.
	\item \textbf{Stacking layers} lets the network build understanding hierarchically: early layers detect edges and textures, middle layers combine these into shapes like eyes or a jawline, and later layers combine those shapes into the concept of a whole face.
	\item \textbf{Fully connected layers}, near the end of the network, take everything that has been detected and use it to make a final decision -- in this case, which of the four categories the face most likely belongs to.
\end{itemize}

\begin{thm}[Feature Map]
The output produced when a filter is slid across an image, highlighting where in the image that filter's pattern was found.
\end{thm}
 
None of these filters are hand-designed by a person. During training, the network starts with random filters and gradually adjusts them, over many examples, until they become useful at telling the categories apart.


\section{How was this model built}

\subsection{What is TensorFlow?}
TensorFlow is an open-source machine learning framework developed by Google. It provides a comprehensive ecosystem for building and deploying machine learning models, including deep learning architectures like CNNs. TensorFlow supports various platforms and allows for efficient computation on CPUs, GPUs, and TPUs.
 
\subsection{What dataset was used to train this model?}
The dataset used to train this model is the "All-Age-Faces Dataset" created by Tsinghua University of Science and Technology, which contains a diverse collection of facial images across different age groups. The dataset includes images of men, women, boys, and girls, allowing the model to learn and classify faces based on age and gender.
 
\subsection{The training process}
The model was trained using TensorFlow/Keras over multiple passes (\emph{epochs}) through the dataset. On each pass, the model made predictions on a batch of training images, measured how far off those predictions were, and nudged its internal filters slightly in the direction that would reduce that error next time.
 
A portion of the dataset was deliberately held back from training, known as the \emph{validation set}, and used only to check how well the model performed on faces it had never seen before. This distinction matters: a model can appear to do very well simply by memorising the training images, without actually learning to generalise. This is known as \emph{overfitting}.
 
\begin{thm}[Overfitting]
When a model becomes very good at the exact examples it was trained on, but performs poorly on new, unseen examples -- similar to a student memorising practice questions instead of understanding the underlying topic.
\end{thm}
 
To help the model generalise better rather than simply memorise, \emph{data augmentation} was applied during training -- automatically creating slightly modified versions of training images (small rotations, flips, zooms) so the model sees more variety without needing more raw data. With this in place, the model reached a validation accuracy of approximately \textbf{77\%} across the four categories.
 
It is worth noting that some categories are inherently harder to separate than others: distinguishing an adult man from a teenage boy, for example, is a genuinely subjective judgement even for a human observer working from a single photo, so some level of confusion between closely related categories is expected rather than a flaw in the model.

\end{document}
